\section{DSC}
\label{sec:DSC}
\subsection{Description des manipulations}

Pour les manipulations d'un echantillon de PET, nous utilisons donc la technique DSC grace a la machine X’Press DSC 214 Polyma' egalement du fournisseur NETZSCH. Cet instrument est essentiellement un four. A l'intereur de ce derniers, se trouvent deux petites capsules visibles sur Fig. \ref{fig:DSCschema}. L'une accueille l'echantillon, l'autre sert de reference. (Ref. \cite{DSC}),

\begin{figure}[ht]
\centering
        \includegraphics[scale=0.55]{schema_dsc.png}
        \captionof{figure}{ Schema de principe de la mesure DSC, Ref. \cite{DSC} }
	\label{fig:DSCschema}
\end{figure}

%	DEFINE CHALEUR SPECIFIQUE + CAPACITE CALO
L'echantillon doit soigneusement etre prepare afin de limite les erreurs de mesures. Il est important de sceller les capsules par pression apres avoir manipuler l'echantillon avec des gants. Cela permet d'eviter d'eviter de contaminer le PET avec des impuretees, car nos mesures se basent sur les capacites calorifiques propres a notre echantillon.
\\\\
%	DEFINE ENTALPIE ...
Enfin, le principe du DSC consiste a simplement mesurer avec deux thermocouples identiques la chaleur que degagent chacun des echantillons. Nous soustrayons la difference entre les deux afin d'obtenir l'evolution de ce parametre en fonction de la temperature du four. En appliquant cette mesure a differentes rampes thermiques (voir vitesses Tab. \ref{val:peaks}), nous pouvons verifier graphiquement la temperature transition vitreuse $V_f$ (ou transition alpha), la temperature de fusion $T_f$. L'experience permet egalement d'identifier les phases ou se forment et deforment les cristaux dans un echantillon de PET en fonction de differentes rampes thermiques. Nous employons l'ordre suivant : d'abord chauffer rapidement l'echantillon, puis le refroidir rapidement. Ensuite il faut proceder a une chauffe lente puis un refroidissement lent. Nous pouvons par consequent extraire le taux de cristalinite a partir de notre mesure d'enthalpie de fusion.

%%%%%%%%%%%%%%%%%%%%%%%%%%%%%%%%%%%%%%%%%%%%%%%
